% GENERATED FILE. Edit canonical lessons, then run npm run generate:books.
% canonical-source-hash: fnv1a64:7945d8683aaaed77
% canonical-lessons: MR-C328-satru, MR-C328-halla, MR-C328-akraman, MR-C328-ladhne, MR-C328-ranangan

\chapter{War Comes}
\label{ch:mr-war-comes}
\hlchaptermodality{328}

\begin{chapteropening}
\textbf{By the end of this chapter:} \emph{I can tell of an enemy, an attack and a battle.}

Complete the last lesson of chapter 328: I can tell of an enemy, an attack and a battle.
\end{chapteropening}

\section[\texorpdfstring{\mr{शत्रू} (śatrū)}{शत्रू (śatrū)}]{\mr{शत्रू} (śatrū) — an enemy}

\label{lesson:MR-C328-satru}



\begin{quote}
Before the new one: say the Marathi for a chieftain, then the Marathi for to found a settlement.
\end{quote}

\subsection*{You'll want to know: \mr{शत्रू}}
\textbf{\mr{शत्रू}} — \emph{śatrū} — \textquotedblleft{}an enemy\textquotedblright{}.

\begin{quote}
\mr{शत्रू} \mr{रात्रीच्या} \mr{अंधारात} \mr{गावात} \mr{शिरला}.
\end{quote}

\begin{quote}
\emph{śatrū rātrīcyā andhārāt gāvāt śirlā.}
\end{quote}

\begin{quote}
The enemy crept into the village in the dark of night.
\end{quote}

Masculine.

Plural \mr{शत्रू}.

\begin{cousinweb}
\textbf{\mr{शत्रू}} (\emph{śatrū}) is a tatsama, taken from Sanskrit \emph{śatru}, 'an enemy, a rival' (Monier-Williams).
\end{cousinweb}

\begin{grammarlens}[title={What you've built}]
One more word for telling what happened.
\end{grammarlens}

\subsection*{Your turn}
\begin{itemize}
\raggedright
  \item \emph{Say it:} \emph{śatrū}
  \item \emph{Say it:} \emph{śatrū}, once more
  \item \emph{Say it:} \emph{śatrū rātrīcyā andhārāt gāvāt śirlā.} — The enemy crept into the village in the dark of night.
  \item \emph{Recall:} say the Marathi for a chieftain, then the Marathi for to found a settlement, then say \emph{śatrū} again
\end{itemize}

\begin{tcolorbox}[breakable,colback=teal!4,colframe=teal!35!black,title={Before you move on}]
What does \mr{शत्रू} mean? (\textquotedblleft{}An enemy\textquotedblright{}.) Say it once more.
\end{tcolorbox}
\section[\texorpdfstring{\mr{हल्ला} (hallā)}{हल्ला (hallā)}]{\mr{हल्ला} (hallā) — an attack}

\label{lesson:MR-C328-halla}



\begin{quote}
Before the new one: say the Marathi for to found a settlement, then the Marathi for an enemy.
\end{quote}

\subsection*{You'll want to know: \mr{हल्ला}}
\textbf{\mr{हल्ला}} — \emph{hallā} — \textquotedblleft{}an attack\textquotedblright{}.

\begin{quote}
\mr{पहाटे} \mr{सैन्यावर} \mr{अचानक} \mr{हल्ला} \mr{झाला}.
\end{quote}

\begin{quote}
\emph{pahāṭe sainyāvar acānak hallā jhālā.}
\end{quote}

\begin{quote}
At dawn the army was suddenly attacked.
\end{quote}

Masculine.

X-\mr{वर} \mr{हल्ला} \mr{करणे}, to attack X.

\begin{cousinweb}
\textbf{\mr{हल्ला}} (\emph{hallā}) is inherited: CDIAL lists Marathi \emph{hallā}, 'an attack', under \emph{*halla}, 'movement' (CDIAL 14017).
\end{cousinweb}

\begin{grammarlens}[title={What you've built}]
One more word for telling what happened.
\end{grammarlens}

\subsection*{Your turn}
\begin{itemize}
\raggedright
  \item \emph{Say it:} \emph{hallā}
  \item \emph{Say it:} \emph{hallā}, once more
  \item \emph{Say it:} \emph{pahāṭe sainyāvar acānak hallā jhālā.} — At dawn the army was suddenly attacked.
  \item \emph{Recall:} say the Marathi for to found a settlement, then the Marathi for an enemy, then say \emph{hallā} again
\end{itemize}

\begin{tcolorbox}[breakable,colback=teal!4,colframe=teal!35!black,title={Before you move on}]
What does \mr{हल्ला} mean? (\textquotedblleft{}An attack\textquotedblright{}.) Say it once more.
\end{tcolorbox}
\section[\texorpdfstring{\mr{आक्रमण} (ākramaṇ)}{आक्रमण (ākramaṇ)}]{\mr{आक्रमण} (ākramaṇ) — an invasion}

\label{lesson:MR-C328-akraman}



\begin{quote}
Before the new one: say the Marathi for an enemy, then the Marathi for an attack.
\end{quote}

\subsection*{You'll want to know: \mr{आक्रमण}}
\textbf{\mr{आक्रमण}} — \emph{ākramaṇ} — \textquotedblleft{}an invasion\textquotedblright{}.

\begin{quote}
\mr{शत्रूच्या} \mr{आक्रमणानंतर} \mr{राज्य} \mr{कमकुवत} \mr{झालं}.
\end{quote}

\begin{quote}
\emph{śatrūcyā ākramaṇānantar rājya kamkuvat jhālaṁ.}
\end{quote}

\begin{quote}
After the enemy's invasion the kingdom grew weak.
\end{quote}

Neuter.

X-\mr{वर} \mr{आक्रमण} \mr{करणे}, to invade X.

\begin{cousinweb}
\textbf{\mr{आक्रमण}} (\emph{ākramaṇ}) is a tatsama, taken from Sanskrit \emph{ākramaṇa}, 'a stepping upon, an attack' (Monier-Williams).
\end{cousinweb}

\begin{grammarlens}[title={What you've built}]
One more word for telling what happened.
\end{grammarlens}

\subsection*{Your turn}
\begin{itemize}
\raggedright
  \item \emph{Say it:} \emph{ākramaṇ}
  \item \emph{Say it:} \emph{ākramaṇ}, once more
  \item \emph{Say it:} \emph{śatrūcyā ākramaṇānantar rājya kamkuvat jhālaṁ.} — After the enemy's invasion the kingdom grew weak.
  \item \emph{Recall:} say the Marathi for an enemy, then the Marathi for an attack, then say \emph{ākramaṇ} again
\end{itemize}

\begin{tcolorbox}[breakable,colback=teal!4,colframe=teal!35!black,title={Before you move on}]
What does \mr{आक्रमण} mean? (\textquotedblleft{}An invasion\textquotedblright{}.) Say it once more.
\end{tcolorbox}
\section[\texorpdfstring{\mr{लढणे} (laḍhṇe)}{लढणे (laḍhṇe)}]{\mr{लढणे} (laḍhṇe) — to fight}

\label{lesson:MR-C328-ladhne}



\begin{quote}
Before the new one: say the Marathi for an attack, then the Marathi for an invasion.
\end{quote}

\subsection*{You'll want to know: \mr{लढणे}}
\textbf{\mr{लढणे}} — \emph{laḍhṇe} — \textquotedblleft{}to fight\textquotedblright{}.

\begin{quote}
\mr{सैनिक} \mr{शेवटपर्यंत} \mr{धैर्याने} \mr{लढले}.
\end{quote}

\begin{quote}
\emph{sainik śevaṭparyant dhairyāne laḍhle.}
\end{quote}

\begin{quote}
The soldiers fought bravely to the end.
\end{quote}

Past (he …): \mr{तो} \mr{लढला}, \emph{to laḍhlā}.

The opponent takes -\mr{शी} or -\mr{विरुद्ध}: \mr{शत्रूशी} \mr{लढणे}.

\begin{cousinweb}
\textbf{\mr{लढणे}} (\emph{laḍhṇe}) is inherited: CDIAL lists Marathi \emph{laḍhṇe}, 'fight', under \emph{*laḍ}, 'fight' (CDIAL 10920).
\end{cousinweb}

\begin{grammarlens}[title={What you've built}]
One more word for telling what happened.
\end{grammarlens}

\subsection*{Your turn}
\begin{itemize}
\raggedright
  \item \emph{Say it:} \emph{laḍhṇe}
  \item \emph{Say it:} \emph{laḍhṇe}, once more
  \item \emph{Say it:} \emph{sainik śevaṭparyant dhairyāne laḍhle.} — The soldiers fought bravely to the end.
  \item \emph{Recall:} say the Marathi for an attack, then the Marathi for an invasion, then say \emph{laḍhṇe} again
\end{itemize}

\begin{tcolorbox}[breakable,colback=teal!4,colframe=teal!35!black,title={Before you move on}]
What does \mr{लढणे} mean? (\textquotedblleft{}To fight\textquotedblright{}.) Say it once more.
\end{tcolorbox}
\section[\texorpdfstring{\mr{रणांगण} (raṇāṅgaṇ)}{रणांगण (raṇāṅgaṇ)}]{\mr{रणांगण} (raṇāṅgaṇ) — a battlefield}

\label{lesson:MR-C328-ranangan}



\begin{quote}
Before the new one: say the Marathi for an invasion, then the Marathi for to fight.
\end{quote}

\subsection*{You'll want to know: \mr{रणांगण}}
\textbf{\mr{रणांगण}} — \emph{raṇāṅgaṇ} — \textquotedblleft{}a battlefield\textquotedblright{}.

\begin{quote}
\mr{रणांगणावर} \mr{दोन्ही} \mr{सैन्यं} \mr{समोरासमोर} \mr{आली}.
\end{quote}

\begin{quote}
\emph{raṇāṅgaṇāvar donhī sainyaṁ samorāsamor ālī.}
\end{quote}

\begin{quote}
On the battlefield the two armies came face to face.
\end{quote}

Neuter.

\mr{रण}, battle, + \mr{अंगण}, a yard.

\begin{cousinweb}
\textbf{\mr{रणांगण}} (\emph{raṇāṅgaṇ}) is built from Sanskrit \emph{raṇa}, 'battle', and \emph{aṅgaṇa}, 'a courtyard' (Monier-Williams): the yard of battle.
\end{cousinweb}

\begin{grammarlens}[title={What you've built}]
One more word for telling what happened.
\end{grammarlens}

\subsection*{Your turn}
\begin{itemize}
\raggedright
  \item \emph{Say it:} \emph{raṇāṅgaṇ}
  \item \emph{Say it:} \emph{raṇāṅgaṇ}, once more
  \item \emph{Say it:} \emph{raṇāṅgaṇāvar donhī sainyaṁ samorāsamor ālī.} — On the battlefield the two armies came face to face.
  \item \emph{Recall:} say the Marathi for an invasion, then the Marathi for to fight, then say \emph{raṇāṅgaṇ} again
\end{itemize}

\begin{tcolorbox}[breakable,colback=teal!4,colframe=teal!35!black,title={Before you move on}]
What does \mr{रणांगण} mean? (\textquotedblleft{}A battlefield\textquotedblright{}.) Say it once more.
\end{tcolorbox}
